\title{DeepSeekMath: Pushing the Limits of Mathematical Reasoning in Open Language Models}

\begin{abstract}
Language models often encounter difficulties with mathematical reasoning because of the domain’s inherent complexity and structure. In this study, we present DeepSeekMath~7B, a model developed by continuing the pre-training of DeepSeek-Coder-Base-v1.5 7B using 120B math-related tokens drawn from Common Crawl, in addition to incorporating both natural language and code resources. DeepSeekMath~7B demonstrates a strong performance by attaining a 51.7\% accuracy on the advanced MATH benchmark, doing so without the aid of auxiliary toolkits or ensembling methods, and reaching near the results of models such as Gemini-Ultra and GPT-4. Furthermore, with self-consistency evaluation across 64 outputs, DeepSeekMath~7B raises its score to 60.9\% on MATH. Two main factors contribute to DeepSeekMath~’s mathematical problem-solving prowess: the use of a carefully crafted pipeline for selecting high-quality public web data and the development of Group Relative Policy Optimization (GRPO), an adaptation of Proximal Policy Optimization (PPO) that not only boosts reasoning performance but also improves PPO’s memory efficiency.
\end{abstract}

\section{Introduction}

Large language models (LLMs) have transformed the field of mathematical reasoning within artificial intelligence, leading to notable progress in benchmarks for quantitative reasoning \citep{MATH} as well as geometry reasoning \citep{trinh2024solving}. In addition, these models have played a pivotal role in supporting humans with challenging mathematical tasks \citep{tao}. Nonetheless, state-of-the-art systems like GPT-4 \citep{gpt4} and Gemini-Ultra \citep{gemini} remain closed-source, with existing open-source alternatives showing a substantial performance gap relative to these leading models.

This work presents DeepSeekMath, a specialized language model designed to excel in mathematical reasoning, achieving superior results compared to open-source alternatives and approaching GPT-4’s proficiency on established academic evaluations. Central to this achievement is the construction of the DeepSeekMath Corpus, a comprehensive and high-quality pre-training resource totaling 120B mathematical tokens. The dataset originates from the Common Crawl (CC) and is filtered using a fastText-based classifier \citep{joulin2016fasttext}. Initially, this classifier is trained with positive samples drawn from OpenWebMath \citep{openwebmath} and a varied assortment of non-mathematical web pages as negative examples. Utilizing this trained classifier, we proceed to identify further positive candidates from the CC, subsequently refining these selections via human annotation. The classifier is then retrained on this augmented set to enhance its discriminative capacity. Empirical results demonstrate the corpus’s strong quality: our foundational DeepSeekMath-Base 7B model attains 64.2\% on GSM8K \citep{gsm8k} and 36.2\% on the MATH competition dataset \citep{MATH}, surpassing the performance of Minerva 540B \citep{minerva}. Additionally, the multilingual nature of the DeepSeekMath Corpus leads to measurable gains on Chinese math benchmarks \citep{wei2023cmath,agieval}. We view our methodology for mathematical data curation as a useful foundation for further investigation, acknowledging considerable opportunity for advancement in this domain.

DeepSeekMath-Base utilizes DeepSeek-Coder-Base-v1.5 7B \citep{deepseek-coder} for initialization, based on our finding that code-oriented pretrained models offer a superior foundation relative to more broadly trained LLMs. In addition, our experiments reveal that mathematical training boosts the model’s performance on both MMLU \citep{mmlu} and BBH \citep{bbh} benchmarks, suggesting that the benefits extend beyond mathematics and encompass general reasoning proficiency as well.

Following the pre-training phase, we conduct mathematical instruction tuning on DeepSeekMath-Base utilizing datasets featuring chain-of-thought \citep{cot}, program-of-thought \citep{pot,pal}, and tool-integrated reasoning \citep{tora}. The trained model, DeepSeekMath-Instruct 7B, surpasses all other 7B models and achieves performance on par with instruction-tuned open-source models at the 70B scale.

Additionally, we present Group Relative Policy Optimization (GRPO), a reinforcement learning (RL) algorithm inspired by Proximal Policy Optimization (PPO) \citep{schulman2017proximal}, but with a crucial modification: GRPO dispenses with the critic model and instead computes the baseline utilizing group scores. This alteration yields a considerable reduction in computational demands during training. Leveraging only a portion of English instruction tuning data, GRPO delivers marked gains over the robust DeepSeekMath-Instruct model, achieving noticeable improvements across both in-domain tasks (GSM8K: 82.9\% $\rightarrow$ 88.2\%, MATH: 46.8\% $\rightarrow$ 51.7\%) and out-of-domain scenarios (e.g., CMATH: 84.6\% $\rightarrow$ 88.8\%) in the reinforcement learning stage. Furthermore, we outline a unified framework to interpret several techniques—namely Rejection Sampling Fine-Tuning (RFT) \citep{yuan2023scaling}, Direct Preference Optimization (DPO) \citep{dpo}, PPO, and GRPO—as either instances of, or simplifications to, conventional RL approaches. Through systematic experimentation—including comparisons of online vs. offline learning, examination of outcome vs. process-based supervision, and analyses contrasting single-turn and iterative RL—we thoroughly explore the critical components underpinning this unified view. Finally, we discuss the mechanisms through which our RL strategy enhances the capabilities of instruction-tuned models and highlight prospective avenues for advancing RL effectiveness within this unifying paradigm.

\subsection{Contributions}
We present scalable approaches to math pre-training and undertake a comprehensive investigation and assessment of reinforcement learning.

\textbf{Math Pre-Training at Scale}

\begin{itemize}[topsep=0pt]
    \item This study demonstrates that Common Crawl's openly available data holds significant potential for mathematical applications. Through a carefully crafted data selection process, we assembled the DeepSeekMath Corpus—a comprehensive and high-caliber dataset comprising 120B tokens extracted from web sources curated for mathematical relevance. This corpus is nearly sevenfold larger than the math-related web data utilized by Minerva \citep{minerva} and approximately nine times greater than the scale of the recently introduced OpenWebMath dataset \citep{openwebmath}.

\item DeepSeekMath-Base 7B, our pre-trained base model, attains performance similar to Minerva 540B \citep{minerva}, suggesting that parameter count alone does not determine effectiveness in mathematical reasoning. High-quality pre-training data can enable a smaller model to deliver competitive results.

\item We present our results from conducting math training experiments. Training models on code before introducing math tasks leads to enhanced performance on mathematical problem solving, regardless of whether external tools are involved. These outcomes contribute to the ongoing debate: \textit{does code training enhance reasoning skills?} Our evidence suggests that it does, specifically in the context of mathematical reasoning.

\item While the practice of utilizing arXiv papers for training is widespread, particularly within the context of mathematics-focused research, it does not yield significant gains on any of the mathematical evaluation benchmarks considered in this study.
\end{itemize}

\textbf{Exploration and Analysis of Reinforcement Learning}

\begin{itemize}[topsep=0pt]
    \item We propose Group Relative Policy Optimization (GRPO), a reinforcement learning technique that prioritizes both efficiency and effectiveness. By omitting the need for a critic model and instead utilizing group scores to compute the baseline, GRPO requires substantially fewer computational resources for training compared to Proximal Policy Optimization (PPO).
    \item We empirically show that applying GRPO notably boosts the capabilities of our instruction-tuned model DeepSeekMath-Instruct, even when relying exclusively on instruction-tuning datasets. In addition, the application of GRPO also yields improved results on tasks that extend beyond the scope of the training domain.
    \item We offer a comprehensive framework that contextualizes and connects various approaches, including RFT, DPO, PPO, and GRPO. Our investigation encompasses an array of experiments—such as comparisons between online and offline learning, between outcome-based and process-based supervision, and between single-turn and multi-step reinforcement learning—aimed at clarifying the foundational components within this framework.
    \item Leveraging this holistic paradigm, we analyze the underlying factors driving the success of reinforcement learning methods and outline prospective research avenues for further improving reinforcement learning of LLMs.
\end{itemize}

\subsection{Summary of Evaluations and Metrics}

\begin{itemize}[topsep=0pt]
    \item \textbf{English and Chinese Mathematical Reasoning}: We perform thorough evaluations of our models using both English and Chinese datasets, which encompass mathematical tasks from elementary through collegiate levels. The English test sets are comprised of GSM8K \citep{gsm8k}, MATH \citep{MATH}, SAT \citep{llemma}, OCW Courses \citep{minerva}, and MMLU-STEM \citep{mmlu}. For Chinese, we make use of benchmarks such as MGSM-zh \citep{mgsm}, CMATH \citep{wei2023cmath}, Gaokao-MathCloze \citep{agieval}, and Gaokao-MathQA \citep{agieval}. Our analysis considers the models' proficiency both in composing independent, explanatory textual solutions (without recourse to tools) and in resolving tasks via Python coding.

When evaluated on English-language benchmarks, DeepSeekMath-Base matches the performance of the proprietary Minerva 540B model \citep{minerva} and consistently outperforms all open-source base models, such as Mistral 7B \citep{mistral} and Llemma-34B \citep{llemma}, regardless of the presence of mathematical pre-training, frequently by a notable margin. Impressively, DeepSeekMath-Base exhibits even greater proficiency on Chinese-language benchmarks, a result that can be attributed to our decision not to limit mathematical pre-training data to English—as was done in previous research \citep{minerva,llemma}—but instead to incorporate high-quality data in additional languages. Following mathematical instruction tuning and reinforcement learning, the subsequent DeepSeekMath-Instruct and DeepSeekMath-RL models deliver robust results, achieving, for the first time among open-source initiatives, accuracy rates exceeding 50\% on the competitive MATH dataset.

\item \textbf{Formal Mathematics}: We assess DeepSeekMath-Base on the informal-to-formal theorem proving benchmark introduced by \citep{dsp_proof}, utilizing the miniF2F dataset \citep{minif2f} with Isabelle \citep{isabelle} as the selected proof assistant. DeepSeekMath-Base attains notable results in few-shot autoformalization.
\item \textbf{Natural Language Understanding, Reasoning, and Code}: To obtain a well-rounded view of the model’s general proficiency in comprehension, logical reasoning, and programming, DeepSeekMath-Base is evaluated using the Massive Multitask Language Understanding (MMLU) benchmark \citep{mmlu}, which includes 57 multiple-choice problems across a wide range of disciplines, the BIG-Bench Hard (BBH) benchmark \citep{bbh}, which comprises 23 complex tasks requiring sophisticated, multi-step inference, as well as widely-adopted code evaluation datasets HumanEval \citep{codex} and MBPP \citep{mbpp}. Math pre-training is shown to improve both the language understanding and reasoning capabilities.

\section{Math Pre-Training}

\subsection{Data Collection and Decontamination} This section details the methodology used to build the DeepSeekMath Corpus by extracting data from Common Crawl. Illustrated in Figure \ref{fig:data_collect}, we introduce a stepwise pipeline designed to methodically assemble an extensive mathematical dataset from Common Crawl, beginning with an initial seed set (for example, a modest yet reliable collection of mathematics-focused data). Importantly, this methodology can be generalized to other areas, such as programming.

Initially, we utilize OpenWebMath \citep{openwebmath}, a curated set of high-quality mathematical web documents, as our foundational seed corpus. With this data, we construct a fastText model \citep{joulin2016fasttext} to retrieve additional web pages that resemble those in OpenWebMath. To achieve this, we draw 500,000 samples at random from the seed corpus as positive examples, while selecting a separate 500,000 web pages from Common Crawl as negative counterparts. The fastText training leverages an open-source implementation\footnote{\url{https://fasttext.cc}}, setting the embedding size to 256, specifying a learning rate of 0.1, capping word n-grams at length 3, imposing a minimum occurrence threshold of 3 for words, and running 3 epochs during training. In order to make the Common Crawl dataset more manageable, we apply both URL-level deduplication and near-duplicate removal techniques, condensing the collection to 40 billion HTML web pages. Next, we employ the trained fastText model to identify mathematical web content from the deduplicated Common Crawl corpus. To eliminate pages of lesser mathematical value, we order the results by the model’s predicted scores, retaining only those with the highest ranks. The extent of data selected is determined by conducting pre-training experiments on subsets consisting of the top 40B, 80B, 120B, and 160B tokens. In our initial cycle, we opt to preserve only the highest-scoring 40B tokens.

Following the initial round of data gathering, a significant number of mathematical web pages are not retrieved, primarily because the fastText model’s positive training set lacks adequate diversity. To address this limitation, we augment the seed corpus by incorporating a broader set of mathematical web sources, thus facilitating enhancements to the fastText model. Our procedure begins with partitioning the entire Common Crawl dataset into non-overlapping domains, where each domain consists of web pages sharing the identical base URL. For every domain, we determine the proportion of pages collected in the initial round. Domains with more than 10\% of their pages retrieved are marked as math-oriented (such as \url{mathoverflow.net}). We then manually label URLs that point to mathematical content within these candidate domains (for instance, \url{mathoverflow.net/questions}). Unretrieved web pages linked to these URLs are subsequently integrated into the seed set. By employing this strategy, we expand the pool of positive training examples, enabling the development of a more effective fastText model with enhanced recall of mathematical data in later cycles. Upon completing four rounds of data collection, we obtain a total of 35.5M mathematical web pages, corresponding to 120B tokens. In the fourth round, it is observed that approximately 98\% of the pages had already been captured during the previous round, leading us to discontinue further data gathering.

In order to prevent benchmark leakage, we adopt the filtering protocol described in \cite{deepseek-coder}, removing web pages that include questions or solutions originating from English mathematical datasets like GSM8K \citep{gsm8k} and MATH \citep{MATH}, as well as Chinese datasets such as CMATH \citep{wei2023cmath} and AGIEval \citep{agieval}. The removal process operates by discarding any passage in our math training data that contains a 10-gram phrase matching exactly with any sequence present in the relevant evaluation benchmarks. For benchmark contents whose length falls below 10 grams but is no less than 3 grams, we apply direct string matching to exclude potentially compromised web pages.

\subsection{Validating the Quality of the DeepSeekMath~Corpus}

We run pre-training experiments to investigate how the DeepSeekMath~Corpus is compared with the recently released math-training corpora:

\begin{itemize}[topsep=0pt]
    \item \textbf{MathPile} \citep{mathpile}: a comprehensive dataset containing 8.9B tokens compiled from diverse sources such as textbooks, Wikipedia, ProofWiki, CommonCrawl, StackExchange, and arXiv; notably, arXiv accounts for more than 85\% of this dataset;
    \item \textbf{OpenWebMath} \citep{openwebmath}: comprises 13.6B tokens of mathematical material selected from CommonCrawl after targeted filtering for mathematical relevance;
    \item \textbf{Proof-Pile-2} \citep{llemma}: this mathematical dataset brings together OpenWebMath, AlgebraicStack (offering 10.3B tokens of math-related code), and scholarly articles from arXiv (28.0B tokens). Consistent with \cite{llemma}, experiments on Proof-Pile-2 employ an arXiv:Web:Code composition ratio of 2:4:1.
\end{itemize}

\subsubsection{Training Setting}
\label{sec:quality-policy}
For our experiments, we start with a general-purpose language model comprising 1.3B parameters, built upon the same architectural foundation as the DeepSeek LLMs \citep{deepseek-llm}; this model is referred to as DeepSeek-LLM 1.3B. Each distinct mathematical dataset is used to train a separate model instance for 150B tokens. We perform all training procedures on the HAI-LLM framework \citep{haillm}, chosen for its speed and efficiency. Adhering to the DeepSeek LLMs training methodology, the optimization employs AdamW \citep{adamW} with hyperparameters set at $\beta_1=0.9$, $\beta_2=0.95$, and $\mathrm{weight\_decay}=0.1$. The learning rate is managed through a multi-step schedule: beginning with a warm-up phase that lasts 2,000 steps to reach the peak rate, it then drops to 31.6\% of its maximum after 80\% of training is complete, and finally decays further to 10.0\% of the maximum at the 90\% training mark. The peak learning rate is established at $5.3\mathrm{e}{-4}$, and training is conducted with a batch size of 4M tokens and a context window of 4K tokens.

\subsubsection{Evaluation Results}

\textbf{The DeepSeekMath~Corpus is of high quality, covers multilingual mathematical content, and is the largest in size.}

\begin{itemize}[topsep=0pt]
    \item \textbf{High-quality}:
    To assess the corpus quality, we tested downstream capabilities across eight mathematical tasks using a few-shot chain-of-thought prompting approach \cite{cot}. According to Table \ref{tab:corpora_comparison}, models trained with the DeepSeekMath~Corpus consistently outperform those trained on alternative datasets. Additionally, Figure \ref{fig:corpora_comparisons} highlights the superior results achieved by the DeepSeekMath Corpus-trained model in comparison to one trained on Proof-Pile-2 at the same token count of 50B (equivalent to one full Proof-Pile-2 epoch), demonstrating the superior average quality present in the DeepSeekMath Corpus.
    \item \textbf{Multilingual}:
    The DeepSeekMath~Corpus contains mathematical content in several languages, with English and Chinese representing the majority of its examples. Table \ref{tab:corpora_comparison} illustrates that models exposed to DeepSeekMath~Corpus gain increased mathematical reasoning abilities in both English and Chinese settings. By contrast, other existing mathematical corpora focus mostly on English, thereby offering less benefit or even detriment to Chinese mathematical reasoning tasks.
    \item \textbf{Large-scale}:
    Compared with prior mathematical corpora, the DeepSeekMath~Corpus provides a much larger dataset. As seen in Figure \ref{fig:corpora_comparisons}, training DeepSeek-LLM 1.3B on DeepSeekMath~Corpus yields a rapidly accelerating learning curve and more persistent advancements over time. Baseline corpora, on the other hand, are substantially smaller and require multiple passes during training, which leads to models quickly hitting a performance ceiling.
\end{itemize}

\subsection{Training and Evaluating DeepSeekMath-Base 7B}

This section presents DeepSeekMath-Base 7B, a foundational model designed to excel at mathematical reasoning. The initial parameters originate from DeepSeek-Coder-Base-v1.5 7B \citep{deepseek-coder}, followed by further pretraining over 500B tokens. The composition of the training dataset is as follows: 56\% sourced from the DeepSeekMath Corpus, 4\% from AlgebraicStack, 10\% from arXiv, 20\% composed of Github code, and the remaining 10\% comprises natural language data extracted from Common Crawl, split between English and Chinese. The training pipeline largely follows the procedures described in Section \ref{sec:quality-policy}, with notable modifications including a peak learning rate of 4.2e-4 and a batch size set at 10M tokens.

This study presents an extensive evaluation of DeepSeekMath-Base 7B’s mathematical proficiency, with particular emphasis on its capacity to generate complete mathematical solutions independently, utilize external tools to address mathematical tasks, and engage in rigorous formal theorem verification. In addition to its mathematical evaluation, we further examine the foundational model’s overall capabilities, covering natural language comprehension, reasoning, and programming abilities.

\paragraph{Mathematical Problem Solving with Step-by-Step Reasoning}

We assess the capabilities of DeepSeekMath-Base in addressing mathematical problems via few-shot chain-of-thought prompting \citep{cot} on eight different benchmarks in both English and Chinese. These datasets include tasks involving quantitative reasoning—such as GSM8K \citep{gsm8k}, MATH \citep{MATH}, and CMATH \citep{wei2023cmath}—as well as multiple-choice question formats, exemplified by MMLU-STEM \citep{mmlu} and Gaokao-MathQA \citep{agieval}. The selected benchmarks span a wide range of mathematical disciplines and difficulty, from elementary mathematics up to problems typical of college-level courses.

Table \ref{tab:base_math_cot} demonstrates that DeepSeekMath-Base 7B achieves superior results across all eight benchmarks when compared to other open-source base models, including the well-established general model Mistral 7B \citep{mistral} and the more recently introduced Llemma 34B \citep{llemma}, which received additional math-specific training using Proof-Pile-2 \citep{llemma}. Remarkably, on the rigorous MATH dataset intended for competition-level tasks, DeepSeekMath-Base outperforms all currently available open-source base models by an absolute margin exceeding 10\%. Furthermore, it achieves better results than Minerva 540B \citep{minerva}—a proprietary model that is 77 times larger, based on PaLM \citep{palm}, and augmented with mathematical literature.

\paragraph{Mathematical Problem Solving with Tool Use} We assess the capability of models to tackle mathematical reasoning tasks from GSM8K and MATH datasets by leveraging few-shot program-of-thought prompting \citep{pot,pal}. For each problem, models are instructed to generate a Python script, permitting the use of libraries like \textit{math} and \textit{sympy} to handle complex calculations. The computed output of the generated code serves as the final solution. According to Table \ref{tab:base_math_pot_proof}, DeepSeekMath-Base 7B demonstrates superior performance relative to the previous best model, Llemma 34B.

\paragraph{Formal Mathematics}

Automating formal proofs offers significant advantages for improving both the correctness and dependability of mathematical arguments, as well as for increasing productivity—a topic that has attracted growing scholarly interest in recent years. In this work, we assess the capabilities of DeepSeekMath-Base 7B on the informal-to-formal proof generation task introduced by \citep{dsp_proof}, which requires constructing a formal proof given an informal problem description, its formal equivalent, and a corresponding informal justification. Our experiments focus on the miniF2F benchmark \citep{minif2f}, designed for evaluating formal reasoning at the level of mathematical Olympiad problems. For each challenge, we prompt the model with a few examples and produce formal proofs in Isabelle. In alignment with the methodology in \cite{dsp_proof}, our approach employs language models to draft outline proofs, after which the general-purpose prover Sledgehammer \citep{sledgehammer} is used to complete any unresolved components. According to the results summarized in Table \ref{tab:base_math_pot_proof}, DeepSeekMath-Base 7B achieves notable success in automating the formalization of proofs.

\paragraph{Natural Language Understanding, Reasoning, and Code}

We assess the natural language understanding capabilities of models using MMLU \citep{mmlu}, their reasoning abilities with BBH \citep{bbh}, and their coding performance through HumanEval \citep{codex} and MBPP \citep{mbpp}. According to the results summarized in Table \ref{tab:understanding_reasoning_code}, DeepSeekMath-Base 7B achieves notable improvements on both MMLU and BBH when compared to its predecessor, DeepSeek-Coder-Base-v1.5 \citep{deepseek-coder}, highlighting the beneficial effects of specialized mathematical training on both comprehension and reasoning tasks. Furthermore, owing to the incorporation of code tokens during continual training, DeepSeekMath-Base 7B successfully preserves the coding performance levels seen in DeepSeek-Coder-Base-v1.5 across both coding benchmarks. In summary, DeepSeekMath-Base 7B markedly surpasses the general-purpose Mistral 7B model \citep{mistral} in the three benchmarks related to reasoning and code generation.

\section{Supervised Fine-Tuning}

\subsection{SFT Data Curation}

We develop a comprehensive instruction-tuning dataset for mathematics that spans both English and Chinese questions, sourced from diverse mathematical domains and encompassing multiple levels of difficulty. Each question is accompanied by corresponding solutions formatted as chain-of-thought (CoT) \citep{cot}, program-of-thought (PoT) \citep{pot,pal}, and tool-integrated reasoning \citep{tora}. In total, the training set consists of 776K examples.

\begin{itemize}[topsep=0pt]
    \item \textbf{English mathematical datasets}: Our resources include GSM8K and MATH, each supplemented with solutions that employ integrated tools. Additionally, we utilize selected examples from MathInstruct \citep{MathInstruct} and the Lila-OOD training partition \citep{lila}, in which questions are addressed via CoT or PoT methods. This compilation represents a broad spectrum of mathematical domains, including, but not limited to, algebra, probability, number theory, calculus, and geometry.
    \item \textbf{Chinese mathematical datasets}: For the Chinese portion, we aggregate K-12 level math problems that cover 76 distinct sub-fields such as linear equations, providing annotated solutions in both CoT style and via tool-integrated reasoning.
\end{itemize}

\subsection{Training and Evaluating DeepSeekMath-Instruct 7B}

This section presents DeepSeekMath-Instruct 7B, a model fine-tuned for mathematical instruction tasks originating from DeepSeekMath-Base. To construct training batches, examples are combined in random order until the total sequence length reaches 4,000 tokens. The training process consists of 500 iterations, using a batch size of 256 and maintaining a fixed learning rate of 5e-5 throughout.

We evaluate models' mathematical performance both without and with tool use, on 4 quantitative reasoning benchmarks in English and Chinese. We benchmark our model against the leading models of the time:

\begin{itemize}[topsep=0pt]
    \item \textbf{Closed-source models} consist of the following: (1) the GPT series, notably GPT-4 \citep{gpt4} and GPT-4 Code Interpreter~\footnote{\url{https://openai.com/blog/chatgpt-plugins\#\#code-interpreter}} which represent the leading versions, (2) Gemini Ultra and Gemini Pro \citep{gemini}, (3) Inflection-2 \citep{inflection-2}, (4) Grok-1~\footnote{\url{https://x.ai/model-card}},
    in addition to several models from China, such as (5) Baichuan-3~\footnote{\url{https://www.baichuan-ai.com}},
    and (6) the most recent version of GLM-4~\footnote{\url{https://open.bigmodel.cn/dev/api\#glm-4}}

    from the GLM family \citep{glm}.

These general-purpose models typically experience multiple alignment stages.

\item \textbf{Open-source models} comprise: general models such as (1) DeepSeek-LLM-Chat 67B \citep{deepseek-llm}, (2) Qwen 72B \citep{qwen}, (3) SeaLLM-v2 7B \citep{seallm}, and (4) ChatGLM3 6B \citep{chatglm3}. In addition, several models are optimized for mathematical tasks, including: (5) InternLM2-Math 20B~\footnote{\url{https://github.com/InternLM/InternLM-Math}}, an extension of InternLM2 that is specifically trained on mathematical data and subsequently undergoes instruction tuning; (6) Math-Shepherd-Mistral 7B, which leverages PPO training \citep{schulman2017proximal} on Mistral 7B \citep{mistral} utilizing a process-supervised reward framework; (7) the WizardMath family \citep{wizardmath}, which enhances Mistral 7B and Llama-2 70B \citep{llama2} models’ mathematical reasoning capabilities through AI-derived instruction tuning (“evolve-instruct”) and PPO, with datasets primarily drawn from GSM8K and MATH; (8) MetaMath 70B \citep{metamath}, where Llama-2 70B is further adapted using an expanded compilation of GSM8K and MATH datasets; (9) ToRA 34B \cite{tora}, a variant of CodeLlama 34B tailored for tool-augmented mathematical reasoning; and (10) MAmmoTH 70B \citep{MathInstruct}, which is Llama-2 70B fine-tuned using instruction data from MathInstruct.

Referencing Table \ref{tab:sft_rl_math}, when tool usage is not permitted, DeepSeekMath-Instruct 7B exhibits robust capabilities in stepwise reasoning. Specifically, on the competitive MATH dataset, our model outperforms every open-source alternative as well as most proprietary systems (including Inflection-2 and Gemini Pro) by a margin of no less than 9\% in absolute terms. This advantage persists even against models with significantly greater parameter counts (such as Qwen 72B) or those that have been specially optimized using math-centric reinforcement learning methods (like WizardMath-v1.1 7B). Although DeepSeekMath-Instruct achieves results on par with Chinese proprietary models such as GLM-4 and Baichuan-3 within the MATH benchmark, it still does not reach the performance levels demonstrated by GPT-4 and Gemini Ultra.

When evaluated in a setting that permits models to combine natural language reasoning with programmatic tool usage for problem solving, DeepSeekMath-Instruct 7B achieves nearly 60\% accuracy on the MATH dataset, outperforming all previously released open-source models. For the remaining benchmarks, our model demonstrates performance on par with DeepSeek-LLM-Chat 67B, the previous state-of-the-art model, despite being one-tenth the size.  
\section{Reinforcement Learning}

\subsection{Group Relative Policy Optimization} Reinforcement learning (RL) has demonstrated significant success in enhancing the mathematical reasoning skills of LLMs beyond the Supervised Fine-Tuning (SFT) phase \citep{wang2023math,wizardmath}. Here, we present Group Relative Policy Optimization (GRPO), our proposed RL method that is both effective and efficient.

\subsubsection{From PPO to GRPO} Proximal Policy Optimization (PPO) \citep{schulman2017proximal} is an actor-critic RL algorithm that is widely used in the RL fine-tuning stage of LLMs \citep{ouyang2022training}. In particular, it optimizes LLMs by maximizing the following surrogate objective:

\begin{equation}
\footnotesize
    \mathcal{J}_{PPO}(\theta) = \mathbb{E}_{q \sim P(Q), o \sim \pi_{\theta_{old}}(O|q)} \frac{1}{|o|} \sum_{t=1}^{|o|} \min \left[ \frac{\pi_\theta(o_{t} | q, o_{<t})}{\pi_{\theta_{old}}(o_{t} | q, o_{<t})} A_{t},\, \text{clip} \left( \frac{\pi_\theta(o_{t} | q, o_{<t})}{\pi_{\theta_{old}}(o_{t} | q, o_{<t})},\, 1 - \varepsilon,\, 1 + \varepsilon \right)  A_{t} \right] ,
\end{equation}

Here, $\pi_{\theta}$ represents the current policy model, while $\pi_{\theta_{old}}$ denotes the previous iteration of the policy. The variables $q$ and $o$ correspond to questions drawn from the dataset and outputs produced by the preceding policy $\pi_{\theta_{old}}$, respectively. The parameter $\varepsilon$ serves as the clipping coefficient used in PPO to enhance training stability. The advantage term, $A_t$, is calculated through Generalized Advantage Estimation (GAE) \citep{gae}, utilizing the future rewards $\{r_{\ge t}\}$ in combination with a value function $V_{\psi}$ that is itself learned. As a result, PPO requires joint optimization of both the policy network and the value function. Furthermore, in order to counteract excessive optimization of the reward model, it is standard practice to incorporate a per-token KL divergence penalty—measured relative to a reference policy—directly into the reward at each token, as proposed in \citep{ouyang2022training}; specifically,

\begin{equation}
    r_{t} = r_\varphi(q, o_{\le t}) - \beta \log\frac{\pi_{\theta}(o_{t}|q, o_{<t})}{\pi_{ref}(o_{t}|q, o_{<t})},
\label{eq:PPO-reward}
\end{equation}

Here, $r_\varphi$ represents the reward model, $\pi_{ref}$ denotes the reference model—commonly set as the original SFT model—and $\beta$ specifies the KL penalty’s weighting factor.

As the value function employed in PPO is typically another model of comparable size as the policy model, it brings a substantial memory and computational burden. Additionally, during RL training, the value function is treated as a baseline in the calculation of the advantage for variance reduction. While in the LLM context, usually only the last token is assigned a reward score by the reward model, which may complicate the training of a value function that is accurate at each token. To address this, as shown in Figure \ref{fig:grpo}, we propose Group Relative Policy Optimization (GRPO), which obviates the need for additional value function approximation as in PPO, and instead uses the average reward of multiple sampled outputs, produced in response to the same question, as the baseline. More specifically, for each question $q$, GRPO samples a group of outputs $\{o_1, o_2, \cdots, o_G\}$  from the old policy  $\pi_{\theta_{old}}$  and then optimizes the policy model by maximizing the following objective:

\begin{equation}
\footnotesize
\begin{split}
    \mathcal{J}_{GRPO}(\theta) &= \mathbb{E}{[q \sim P(Q), \{o_i\}_{i=1}^G \sim \pi_{\theta_{old}}(O|q)]}  \\
    & \frac{1}{G}\sum_{i=1}^G\frac{1}{|o_i|} \sum_{t=1}^{|o_i|} \left\{ \min \left[ \frac{\pi_\theta(o_{i,t} | q, o_{i,<t})}{\pi_{\theta_{old}}(o_{i,t} | q, o_{i,<t})} \hat{A}_{i,t}, \text{clip} \left( \frac{\pi_\theta(o_{i,t} | q, o_{i,<t})}{\pi_{\theta_{old}}(o_{i,t} | q, o_{i,<t})}, 1 - \varepsilon, 1 + \varepsilon \right)  \hat{A}_{i,t} \right] - \beta \mathbb{D}_{KL}\left[\pi_{\theta} || \pi_{ref}\right]\right\} ,
\end{split}
\label{eq:GRPO-obj}
\end{equation}

The equation above expresses the objective function for GRPO as follows: we compute the expectation over queries $q$ sampled from $P(Q)$ and response sets $\{o_i\}_{i=1}^G$ drawn from the prior policy $\pi_{\theta_{old}}(O|q)$. The objective consists of a double sum—first averaging over $G$ sampled outputs and then over each time-step $t$ in every sampled output $o_i$. At each time-step, the term inside the summation is the minimum between the standard policy ratio (multiplied by the estimated advantage $\hat{A}_{i,t}$) and its value after being clipped within the range $[1-\varepsilon, 1+\varepsilon]$—this is designed to prevent excessively large policy updates. Lastly, the term subtracts a scaled Kullback-Leibler divergence between the current policy $\pi_\theta$ and a reference policy $\pi_{ref}$, weighted by the factor $\beta$.

where $\varepsilon$ and $\beta$ are hyper-parameters, and $\hat{A}_{i,t}$  is the advantage calculated based on relative rewards of the outputs inside each group only, which will be detailed in the following subsections. The group relative way that GRPO leverages to calculate the advantages, aligns well with the comparative nature of rewards models,  as reward models are typically trained on datasets of comparisons between outputs on the same question. Also note that, instead of adding KL penalty in the reward, GRPO regularizes by directly adding the KL divergence between the trained policy and the reference policy to the loss, avoiding complicating the calculation of  $\hat{A}_{i,t}$. And different from the KL penalty term used in (\ref{eq:PPO-reward}), we estimate the KL divergence with the following unbiased estimator \citep{kl_approx}:

\begin{equation}
\small
    \mathbb{D}_{KL}\left[\pi_{\theta} || \pi_{ref}\right] = \frac{\pi_{ref}(o_{i,t}|q,o_{i,<t})}{\pi_{\theta}(o_{i,t}|q,o_{i,<t})} - \log\frac{\pi_{ref}(o_{i,t}|q,o_{i,<t})}{\pi_{\theta}(o_{i,t}|q,o_{i,<t})} - 1,
\end{equation}

which is guaranteed to be positive.

\begin{algorithm}[t]
  \small
  \caption{Iterative Group Relative Policy Optimization}
  \textbf{Input} initial policy network $\pi_{\theta_{\text{init}}}$; reward evaluators $r_\varphi$; set of task prompts $\mathcal{D}$; 
  hyperparameters $\varepsilon$, $\beta$, $\mu$
  \begin{algorithmic}[1]
    \State Initialize policy network: $\pi_\theta \gets \pi_{\theta_{\text{init}}}$
    \For{iteration $= 1$ to $I$}
       \State Set reference model: $\pi_{ref} \gets \pi_{\theta}$
      \For{step $= 1$ to $M$}
      \State Draw a mini-batch $\mathcal{D}_b$ from $\mathcal{D}$
      \State Assign $\pi_{\theta_{old}} \gets \pi_\theta$ to save previous parameters
      \State For each prompt $q \in \mathcal{D}_b$, sample $G$ candidate responses $\{o_i\}_{i=1}^G \sim \pi_{\theta_{old}} (\cdot \mid q)$
      \State For each $o_i$, obtain reward values $\{r_i\}_{i=1}^{G}$ by querying $r_\varphi$
      \State For each token $t$ of $o_i$, estimate group relative advantages $\hat{A}_{i,t}$
      \For{GRPO iteration $= 1$ to $\mu$}
        \State Optimize $\pi_\theta$ by maximizing the group-relative objective in Equation \ref{eq:GC-GRPO}
      \EndFor
    \EndFor 
    \State Refine $r_\varphi$ with ongoing training using a replay buffer mechanism. 
    \EndFor 
  \end{algorithmic}
  \textbf{Output} $\pi_\theta$
  \label{alg:iter-grpo}
\end{algorithm}

\subsubsection{Outcome Supervision RL with GRPO} Specifically, given a question $q$, a set of responses $\{o_1, o_2, \cdots, o_G\}$ is generated by sampling from the previous iteration of the policy, $\pi_{\theta_{old}}$. Each response is evaluated using a reward model, producing a set of $G$ reward values, $\mathbf{r} = \{r_1, r_2, \cdots, r_G\}$. These rewards are then standardized within the group by removing the mean and scaling by the standard deviation. Under outcome supervision, the normalized reward associated with each output $o_i$ is assigned uniformly to all token advantages $\hat{A}_{i, t}$ in that output, expressed as $\hat{A}_{i, t} = \widetilde{r}_i = \frac{r_i- {\rm mean}(\mathbf{r})}{{\rm std}(\mathbf{r})}$. The policy is subsequently updated by maximizing the objective provided in equation (\ref{eq:GRPO-obj}).

\subsubsection{Process Supervision RL with GRPO} While outcome supervision relies on providing feedback solely at the end of each complete output, this may not suffice for effectively guiding the policy when tackling intricate mathematical problems. Inspired by \cite{wang2023math}, our approach incorporates process supervision, whereby feedback is furnished at the conclusion of every individual reasoning step. To describe this formally, consider the question $q$ alongside a set of $G$ sampled responses $\{o_1, o_2, \cdots, o_G\}$. For each step in these responses, a process reward model evaluates and assigns scores, producing reward collections as follows: $\mathbf{R} = \{ \{r_1^{index(1)},\cdots,r_1^{index(K_1)}\}, \cdots,  \{r_G^{index(1)},\cdots,r_G^{index(K_G)}\} \}$, where $index(j)$ specifies the end token index of the $j$-th step in a sample and $K_i$ indicates the total count of reasoning steps within the $i$-th output. These rewards are standardized across all samples using their global mean and standard deviation: $\widetilde{r}_i^{index(j)} = \frac{r_i^{index(j)} - {\rm mean(\mathbf{R})}}{{\rm std(\mathbf{R})}}$.

Following this, for each token, the process supervision framework computes its advantage by summing the normalized rewards associated with all subsequent steps, as formulated by $\hat{A}_{i, t} = \sum_{index(j) \ge t} \widetilde{r}_i^{index(j)}$.
The policy is then improved by maximizing the objective stated in equation (\ref{eq:GRPO-obj}).

\subsubsection{Iterative RL with GRPO} During the reinforcement learning process, the previously trained reward model may become inadequate for guiding the updated policy model. To address this, we investigate an iterative approach to RL with GRPO. As detailed in Algorithm \ref{alg:iter-grpo}, this method involves periodically generating new datasets for the reward model from samples obtained using the current policy model. The reward model is further updated using a replay mechanism that retains 10\% of historical samples. Subsequently, the reference model is updated to match the policy model, and the policy model continues training with the refreshed reward model.

\subsection{Training and Evaluating DeepSeekMath-RL}

Our RL experiments utilize DeepSeekMath-Instruct 7B as the foundation. The RL training set comprises approximately 144K chain-of-thought-style questions derived from GSM8K and MATH subsets within the SFT dataset. We deliberately omit SFT samples from other sources in order to assess how RL influences benchmarks that are not represented during this phase. The construction process for the reward model training set aligns with the methodology in \citep{wang2023math}. For initializing the reward model, we employ DeepSeekMath-Base 7B, adopting a learning rate of 2e-5. Within GRPO training, the policy model uses a learning rate of 1e-6, and the KL penalty is set at 0.04. Each question is used to generate $64$ sampled responses. We constrain both the maximum sequence length and the training batch size to 1024 tokens and samples, respectively. The policy model undergoes only one update for each exploration iteration. When assessing DeepSeekMath-RL 7B, we use the same benchmarks as for DeepSeekMath-Instruct 7B. For evaluation purposes, tasks involving chain-of-thought reasoning from GSM8K and MATH are categorized as in-domain, while all other test sets are treated as out-of-domain.

Table \ref{tab:sft_rl_math} presents a comparison of open- and closed-source models, evaluating their abilities with both chain-of-thought and tool-integrated reasoning approaches on English and Chinese assessment sets. Our observations include: 1) DeepSeekMath-RL 7B achieves accuracy rates of 88.2\% on GSM8K and 51.7\% on MATH when employing chain-of-thought reasoning. These results exceed those of all open-source competitors within the 7B–70B parameter spectrum, as well as most closed-source alternatives. 2) Notably, DeepSeekMath-RL 7B’s training was restricted to instruction tuning data in chain-of-thought format from GSM8K and MATH, initialized from DeepSeekMath-Instruct 7B. In spite of the limited variety in its training inputs, DeepSeekMath-RL 7B delivers superior outcomes over DeepSeekMath-Instruct 7B in all evaluation criteria, underscoring the advantages gained through reinforcement learning.

\section{Discussion}
This section presents the outcomes observed during our pre-training and RL experimental phases.

\subsection{Lessons Learnt in Pre-Training}

We begin by detailing our pre-training process. Throughout, unless explicitly noted, the experimental setup follows the configuration described in Section \ref{sec:quality-policy}. Importantly, for all mentions of the DeepSeekMath Corpus within this discussion, we are referring to an 89-billion-token subset derived from the dataset’s second collection phase.

\subsubsection{Code Training Benefits Mathematical Reasoning}

An often-cited but unproven theory proposes that code training enhances reasoning skills. In this work, we seek to provide preliminary evidence for this claim specifically in the realm of mathematics: code training leads to better mathematical reasoning performance by models, regardless of whether external tools are employed.

To study how code training affects mathematical reasoning, we experimented with the following two-stage training and one-stage training settings:

\textbf{Two-Stage Training}

\begin{itemize}[topsep=0pt]
    \item \textbf{Code Training for 400B Tokens $\rightarrow$ Math Training for 150B Tokens}: The DeepSeek-LLM 1.3B model is initially pre-trained on 400B tokens consisting of code, after which it is further trained on 150B tokens containing mathematical content;
    \item \textbf{General Training for 400B Tokens $\rightarrow$ Math Training for 150B Tokens}: To serve as a baseline, we replace code tokens with general tokens (randomly drawn from a broad, large-scale corpus assembled by DeepSeek-AI) during the initial 400B-token pre-training phase. This design aims to analyze whether code pre-training yields greater improvements in mathematical reasoning than generic data.
\end{itemize}

\textbf{One-Stage Training}

\begin{itemize}[topsep=0pt]
    \item \textbf{150B Math Token Training}: DeepSeek-LLM 1.3B is trained using 150B tokens specifically from mathematical data;
    \item \textbf{Combined Training with 400B Code Tokens and 150B Math Tokens}: Incorporating math training after code training tends to lower coding ability. We therefore explore if a one-stage training process, where code and math tokens are jointly presented, can enhance mathematical reasoning capabilities without causing significant degradation in coding performance, thereby addressing the issue of catastrophic forgetting.
\end{itemize}

\paragraph{Results}
The downstream outcomes for various training configurations are presented in Table \ref{tab:code-math} and Table \ref{tab:code-math-general-eval}.

Training with code contributes positively to program-aided mathematical reasoning, regardless of whether a two-stage or one-stage training approach is employed. Table \ref{tab:code-math} demonstrates that, in the context of two-stage training, exposure to code alone already substantially improves performance on GSM8K and MATH problem-solving tasks when using Python. Introducing math-focused training in the subsequent stage leads to even greater gains. Notably, in the one-stage training scenario, the combination of code tokens and math tokens effectively addresses the catastrophic forgetting typically associated with the two-stage framework and further enhances both coding skills (Table \ref{tab:code-math-general-eval}) and program-aided mathematical reasoning capabilities (Table \ref{tab:code-math}).

Training on code data also enhances mathematical reasoning even when tools are not employed. In the two-stage training paradigm, the preliminary phase focused on code yields noticeable improvements. Furthermore, it facilitates greater effectiveness during subsequent mathematics-specific training, ultimately achieving superior overall results. Conversely, training on a mixture of code and mathematical tokens in a single stage appears to undermine mathematical reasoning in the absence of tool use. A possible explanation is that DeepSeek-LLM 1.3B, given its smaller model size, does not possess sufficient capacity to concurrently learn from both code and mathematical content.

\subsubsection{ArXiv Papers Seem Ineffective in Improving Mathematical Reasoning}

ArXiv papers are commonly included as a component of math pre-training data \citep{minerva,gpt-f,llemma,mathpile}. However, detailed analysis regarding their impact on mathematical reasoning has not been extensively conducted. Perhaps counter-intuitively, according to our experiments, arXiv papers seem ineffective in improving mathematical reasoning. We experiment with models of different sizes, including DeepSeek-LLM 1.3B and DeepSeek-Coder-Base-v1.5 7B \citep{deepseek-coder}, using arXiv corpora that underwent varied processing pipelines:

\begin{itemize}[topsep=0pt]
    \item \textbf{MathPile} \citep{mathpile}: contains 8.9B tokens and was assembled using specific cleaning and filtering heuristics; notably, more than 85\% of its content comprises scientific articles sourced from arXiv.
    \item \textbf{ArXiv-RedPajama} \citep{redpajama}: features 28.0B tokens, consisting of all arXiv LaTeX source files with extraneous content such as preambles, comments, macros, and bibliography entries stripped away.
\end{itemize}

Within our study, we conduct separate training sessions for DeepSeek-LLM 1.3B on 150B tokens and DeepSeek-Coder-Base-v1.5 7B on 40B tokens, utilizing only the respective arXiv corpora. The data indicate that arXiv papers do not significantly enhance mathematical reasoning capabilities. When exposed solely to arXiv content, neither model exhibits any meaningful improvement— and in some cases, even a decline—across the spectrum of mathematical benchmarks of varying difficulty adopted in this research. These evaluation sets include quantitative reasoning tasks such as GSM8K and MATH (see Table \ref{tab:arxiv-cot}), multiple-choice problems exemplified by MMLU-STEM (see Table \ref{tab:arxiv-cot}), as well as formal mathematics assessments like miniF2F (refer to Table \ref{tab:arxiv_atp}).

However, this conclusion has its limitations and should be taken with a grain of salt. We have not yet studied:

\begin{itemize}[topsep=0pt]
    \item How arXiv tokens influence math-related applications not examined in this study, including tasks like informalization of theorems, where formal statements or proofs are translated into informal language;
    \item The interaction between arXiv tokens and additional data sources;
    \item If advantages derived from arXiv papers persist or amplify as model size increases.
\end{itemize}

Thus, further exploration is required, which we leave for future studies.

\subsection{Insights of Reinforcement Learning}
\subsubsection{Towards to a Unified Paradigm} In this part, we introduce a comprehensive framework for examining a range of training techniques, including SFT, RFT, DPO, PPO, and GRPO. We also carry out empirical studies to investigate the key components within this unified framework. Typically, the gradient with respect to the parameter $\theta$ for any such training approach can be expressed as:

\begin{equation}
    \nabla_{\theta}\mathcal{J}_{\textcolor{red}{\mathcal{A}}}(\theta) = \mathbb{E}[\underbrace{(q,o) \sim \textcolor{red}{\mathcal{D}}}_{Data \ Source}]\left( \frac{1}{|o|} \sum_{t=1}^{|o|}  \underbrace{GC_{{\mathcal{A}}}(q, o, t, \textcolor{red}{\pi_{{rf}}})}_{Gradient \ Coefficient}  \nabla_{\theta}\log \pi_{\theta}(o_t | q, o_{<t})\right).
\label{eq:objective}
\end{equation}

There exist three key components: 1) \textit{Data Source $\mathcal{D}$}, which determines the training data; 2) \textit{Reward Function $\pi_{{rf}}$}, which is the source of the training reward signal; 3) \textit{Algorithm $\mathcal{A}$}: which processes the training data and the reward signal to the gradient coefficient $GC$ that determines the magnitude of the penalty or reinforcement for the data. We analyze several representative methods based on such a unified paradigm:

\begin{itemize}
    \item  \textbf{Supervised Fine-tuning (SFT)}: In SFT, the pretrained model undergoes further training on datasets curated and selected by humans.
    \item \textbf{Rejection Sampling Fine-tuning (RFT)}: RFT involves additional fine-tuning of the SFT-trained model by utilizing outputs generated from SFT model responses to SFT prompts, retaining only those responses whose answers are deemed correct for further training.
    \item \textbf{Direct Preference Optimization (DPO)}: DPO builds upon the SFT model by performing additional fine-tuning on augmented samples generated from the SFT model, leveraging a pairwise DPO loss function to guide optimization.
    \item \textbf{Online Rejection Sampling Fine-tuning (Online RFT)}: Unlike RFT, Online RFT begins by initializing the policy model with the SFT-trained parameters and iteratively improves it via fine-tuning on augmented outputs obtained from the current, real-time policy model.
    \item \textbf{PPO/GRPO}: PPO/GRPO starts from the SFT-pretrained policy model, strengthening it with outputs drawn from the current policy model in an online manner.
\end{itemize}

Table \ref{tab:data_policy} presents an overview of the key elements comprising these methods. For a comprehensive explanation of the derivation procedure, consult Appendix \ref{app:analysis-rl}.

\paragraph{Observation about Data Source} We categorize data sources into two types: online sampling and offline sampling. Online sampling refers to training data generated through the exploratory outputs of the currently training policy model, whereas offline sampling uses data derived from the sampling results of the original SFT model. RFT and DPO employ the offline approach, while Online RFT and GRPO utilize online sampling.

As illustrated in Figure \ref{fig:rl-analysis}, our results indicate that Online RFT consistently surpasses RFT across both evaluated benchmarks. In the initial phases of training, the performance of Online RFT and RFT is nearly indistinguishable; however, as training progresses, Online RFT achieves a clear and substantial lead, highlighting the effectiveness of online learning strategies. This observation aligns with expectations: during the early steps, the actor and SFT model are highly similar, resulting in sampled data that differ only marginally. As training advances, discrepancies between the actor's outputs and the SFT model grow, making the benefits of online, real-time data sampling increasingly pronounced.

\paragraph{Observation about Gradient Coefficient} The algorithm leverages the reward signal through the gradient coefficient to facilitate updates to the model parameters. In our study, we decompose the reward function into two categories: `Rule' and `Model.' The `Rule' category involves evaluating the response quality through direct correctness assessment, while the `Model' category refers to employing a trained reward model that assigns scores to responses. The supervision for the reward model originates from the outcomes of rule-based evaluation. As detailed in Equations \ref{eq:GC-RFT} and \ref{eq:GC-GRPO}, there exists a fundamental distinction between GRPO and Online RFT: GRPO distinctively tailors its gradient coefficient according to the magnitude of the reward model's output. This mechanism enables responses to receive varied degrees of reinforcement or penalization proportional to their individual scores. Conversely, Online RFT does not possess this capability; it neither penalizes incorrect answers nor differentiates the degree of reinforcement, instead providing the same positive update to all correct responses irrespective of their specific quality.

Figure \ref{fig:rl-analysis} illustrates that GRPO outperforms online RFT, underscoring the effectiveness of modifying the coefficients for positive and negative gradients. Moreover, GRPO+PS achieves better outcomes than GRPO+OS, suggesting that employing fine-grained, step-dependent gradient coefficients is advantageous. Additionally, we investigate iterative RL by performing two iterations in our experiments. The results in Figure \ref{fig:iter-rl} reveal that iterative RL notably boosts performance, with the most substantial gains observed during the initial iteration.

\subsubsection{Why RL Works?} In this study, we apply reinforcement learning to a selected portion of instruction tuning data, resulting in notable gains over the baseline instruction tuning approach. To investigate the underlying mechanisms of this improvement, we assess both Pass@K and Maj@K accuracies for the Instruct and RL models across two evaluation benchmarks. The results, depicted in Figure \ref{fig:maj-pass}, reveal that reinforcement learning improves Maj@K accuracy but does not affect Pass@K. This suggests that RL primarily strengthens the overall distribution of outputs, \textbf{improving the likelihood that the correct answer appears among the top K candidates, rather than fundamentally advancing the model’s core abilities.} In a similar vein, \citep{wang2023making} reported a \textbf{misalignment problem} with reasoning tasks in SFT models, demonstrating that preference alignment methods can effectively improve the reasoning performance of these models \citep{yuan2023rrhf,song2023preference,wang2023making}.

\subsubsection{How to Achieve More Effective RL?}
\label{sec:effec_rl}
Our findings indicate that RL performs remarkably well on mathematical reasoning problems. Additionally, we introduce a cohesive framework that encapsulates various notable training strategies, interpreting them either as explicit applications of RL or as approximations thereof. According to the structure outlined in Equation \ref{eq:objective}, three fundamental elements emerge: Data Source, Algorithm, and Reward Function. Below, we outline several prospective avenues for advancing each of these components.

\paragraph{Data Source} The data source forms the foundational input for all training approaches. Within the scope of RL, this term denotes unlabeled questions paired with outputs generated from the policy model. For this study, we limit our data to questions derived from the instruction tuning phase, utilizing a straightforward nucleus sampling technique to obtain responses. This limitation may contribute to our RL pipeline’s improvement being observed only in Maj@K metrics. Looking ahead, we intend to apply our RL pipeline to prompts that fall outside the training distribution and to integrate \textbf{advanced sampling (decoding) strategies}, including those utilizing tree-search algorithms \citep{yao2023tree}. Furthermore, the adoption of \textbf{efficient inference techniques} \citep{xia-etal-2023-speculative,leviathan2023fast,kwon2023efficient,xia2024unlocking}, which are critical for optimizing the exploration process of policy models, is recognized as being highly significant.

\paragraph{Algorithms} Algorithms utilize both the data and the reward signal in determining the gradient coefficient, which in turn is used for model parameter updates. According to Equation \ref{eq:objective}, existing approaches essentially place full \textbf{TRUST} in the information provided by the reward function, adjusting the conditional probability of specific tokens upwards or downwards based on this signal. Nonetheless, the reliability of the reward signal cannot be universally guaranteed, particularly when dealing with tasks of high complexity. As an illustration, even in meticulously annotated datasets like PRM800K \citep{lightman2023let}, produced by highly trained annotators, an estimated 20\% of the annotations are incorrect\footnote{\url{https://github.com/openai/prm800k/issues/12\#issuecomment-1728491852}}. In light of these issues, we intend to investigate reinforcement learning algorithms that maintain robustness in the face of noisy reward signals. We posit that adopting \textbf{WEAK-TO-STRONG} \citep{burns2023weak} alignment strategies may fundamentally transform how learning algorithms operate.

\paragraph{Reward Function} The reward function serves as the primary training signal. In reinforcement learning, this function is commonly implemented through a neural reward model. We identify three crucial avenues for advancing reward models: 1) \textbf{Improving the generalization capacity of the reward model.} It is essential for the reward model to generalize well to out-of-distribution queries and sophisticated decoding outputs; without such robustness, reinforcement learning might simply reinforce the existing LLM distribution without fostering true progress in capability; 2) \textbf{Incorporating uncertainty estimation within the reward model.} The assessment of uncertainty may facilitate more effective connections between less robust reward models and learning algorithms that transition from weak to strong performance; 3) \textbf{Developing efficient methodologies for constructing high-quality process reward models} capable of delivering nuanced, detailed feedback during the reasoning process \citep{lightman2023let,wang2023math}.

\section{Conclusion, Limitation, and Future Work}
This paper introduces DeepSeekMath, which achieves superior results compared to all other open-source models on the MATH competition-level benchmark, and comes close to matching the scores obtained by proprietary models. The model leverages DeepSeek-Coder-v1.5 7B as its foundation and is subject to an extended continual pre-training phase over 500B tokens, including a substantial portion of 120B mathematics-focused tokens extracted from Common Crawl. Through a comprehensive ablation analysis, our findings indicate that web sources provide abundant, high-quality mathematical datasets, whereas arXiv does not deliver the anticipated advantages. We further propose Group Relative Policy Optimization (GRPO)—a modification of Proximal Policy Optimization (PPO)—demonstrating it substantially strengthens mathematical reasoning performance while maintaining lower memory usage. Experimental outcomes highlight GRPO’s effectiveness, as shown by the robust benchmark performance of DeepSeekMath-Instruct 7B. Additionally, we outline a cohesive framework for interpreting a collection of existing approaches and discuss prospective research avenues to improve reinforcement learning methodologies.

While DeepSeekMath~demonstrates strong performance on various quantitative reasoning tasks, its abilities in geometry and theorem proving lag behind those of proprietary models. For example, our preliminary evaluation reveals that the model struggles with questions involving triangles and ellipses, suggesting a potential bias in the pre-training and fine-tuning data used. Moreover, due to its relatively limited model size, DeepSeekMath~underperforms GPT-4 in few-shot learning scenarios. Whereas GPT-4 exhibits notable improvement when provided with few-shot examples, DeepSeekMath~yields comparable results under both zero-shot and few-shot settings. Looking ahead, we plan to refine our data selection processes to build a higher-quality pre-training dataset. We also intend to investigate new approaches for enhancing the effectiveness of reinforcement learning in large language models, as discussed in Section \ref{sec:effec_rl}.

\end{document}